\documentclass[12pt,a4paper]{exam}

% --------------------
% Packages
% --------------------
\usepackage[utf8]{inputenc}
\usepackage[T1]{fontenc}
\usepackage[french]{babel}
\usepackage{amsmath,amssymb}
\usepackage{multicol}
\usepackage{geometry}
\usepackage{pgfplots}

\pgfplotsset{compat=1.18}
\geometry{margin=2cm}

% --------------------
% En-tête / pied (exam)
% --------------------
\header{}{Interrogation de mathématiques}{2nde10}
\footer{}{\thepage}{\it{G.M.}}

% --------------------
\begin{document}

\begin{center}
    \Large\textbf{Interrogation\\Équations et inéquations}
\end{center}

\vspace{0.3cm}

\noindent
\textbf{Durée :} 1h \hfill
\textbf{Calculatrice :} \underline{NON} autorisée

\vspace{0.5cm}

% ======================================================
\noindent
\textbf{Exercice 1 : Équations (\ \ \ \  / 9 points )}
\vspace{5pt}

\begin{questions}

\question[1] Résoudre les équations suivantes dans $\mathbb{R}$ :
\begin{multicols}{2}
    \begin{parts}
        \part $4x + 3 = 3x$
        \part $\dfrac{4}{3} -\dfrac{4}{5}x = 0$
    \end{parts}
\end{multicols}
\begin{multicols}{2}
    \begin{parts}
        \setcounter{partno}{2}
        \part $x^2 - 5 = 0$
        \part $(4x + 2)(3x - 4) = 0$
    \end{parts}
\end{multicols}
\begin{multicols}{2}
    \begin{parts}
        \setcounter{partno}{4}
        \part $\dfrac{2x - 1}{x+7} = 0$
        \part $\dfrac{x(3x+2)}{x^2+3} = 0$
    \end{parts}
\end{multicols}

\end{questions}

\vspace{0.5cm}

% ======================================================
\noindent
\textbf{Exercice 2 : Inéquations (\ \ \ \  / 9 points )}
\vspace{5pt}

\begin{questions}

\question[4] Résoudre les inéquations suivantes dans $\mathbb{R}$ et indiquer les solutions 
\\\textbf{\underline{sous forme d'intervalles}} : $S = ...$
\begin{multicols}{2}
    \begin{parts}
        \part $8x + 3 > 3x - 5$
        \part $-3x \geq 3x - 1$
    \end{parts}
\end{multicols}
\begin{multicols}{2}
    \begin{parts}
        \setcounter{partno}{2}
        \part $(1-2x)(4x+2) \leq 0$
        \part $\dfrac{-8x+16}{2x} < 0$
    \end{parts}
\end{multicols}

\end{questions}

\vspace{10pt}

% ======================================================
\noindent
\textbf{Exercice 3 : Approfondissement (\ \ \ \  / 2 points )}
\vspace{5pt}

\begin{questions}

\question[1] On considère la fonction $g$ suivante :
\[
g(x) = \dfrac{1}{x} - \dfrac{2}{5-x} \quad \text{pour tout $x \neq 0$ et $x \neq 5$}
\]
Résoudre $g(x) \leq 0$.

\vspace{5pt}

\question[1] Résoudre dans $\mathbb{R}$ :
\begin{multicols}{2}
    \begin{parts}
        \part $(1-2x)(3x+1) > (1-x)(3+9x)$
        \part $(11x-3)^2 = 11$
    \end{parts}
\end{multicols}

\end{questions}

\vspace{10pt}

% ======================================================
\noindent
\textbf{Exercice 4 : BONUS (\ \ \ \  / 2 points )}
\vspace{5pt}

\begin{questions}

\question Résoudre dans $\mathbb{R}$ :
\[
\dfrac{9x^2-42x+49}{9x^2-49} = 0
\]

\end{questions}

\vspace{10pt}

\begin{center}
    \textit{Fin du devoir}
\end{center}

\end{document}
